% PROBLEMI CON CUI DOVREMMO FARE LE RIDUZIONI ALL'ESAME
\chapter{Problemi difficili e completi*}
\section{Introduzione}
Vediamo in cosa consistono i \textbf{problemi difficili e completi}:\smallskip\\
\colorbox{SALZO}{\parbox{\dimexpr\linewidth-2\fboxsep}{\begin{definition}
   Un problema $\mathcal{P}$ è $\bm{\mathbb{C}-\textbf{hard}}$ se una Macchina di Turing che risolve $\mathcal{P}$
   può risolvere ogni problema nella clasee $\mathbb{C}$ se utilizzata come subroutine di da altre macchine
\end{definition}}}\smallskip\\
In sostanza vuol dire che il problema $\mathcal{P}$ è almeno tanto difficile quanto ogni altro problema 
in $\mathbb{C}$, se non di più.\smallskip\\
\colorbox{SALZO}{\parbox{\dimexpr\linewidth-2\fboxsep}{\begin{definition}
   Un problema $\mathcal{P}$ è $\bm{\mathbf{\mathbb{C}}-\textbf{complete}}$ se è $\mathbb{C}$-hard e appartiene a $\mathbb{C}$
\end{definition}}}\smallskip\\
Qui invece significa che $\mathcal{P}$ è tanto difficile quanto ogni altro problema in $\mathbb{C}$, ma è anche in $\mathbb{C}$.\\
Una macchina per $\mathcal{P}$ risolve ogni problema in $\mathbb{C}$ se, per ogni problema $\mathcal{P'}$ di 
$\mathbb{C}$ esiste una \textbf{riduzione} da $\mathcal{P}$ a $\mathcal{P'}$.\\
Ricordiamo che $A$ è riducibile (molti a uno) a $B$ ($A \le_m B$), se esiste una funzione calcolabile 
$f:\Sigma_{A}^* \to \Sigma_B^*$ tale che, per ogni $x \in \Sigma_A^*$, $x \in A$ se e solo se $f(x) \in B$. 
(La spiegazione completa delle riduzioni sta qui \ref{riduzioni}).\smallskip\\
Utilizziamo il concetto di riduzione per definire i \textbf{linguaggi difficili} ($\mathbb{C}$-hard):\smallskip\\
\colorbox{SALZO}{\parbox{\dimexpr\linewidth-2\fboxsep}{\begin{definition}
    Un linguaggio $\mathcal{L}$ è $\bm{\mathbb{C}-\textbf{hard}}$ se, per ogni $\mathcal{L}' \in \mathbb{C}$ 
    esiste una macchina di Turing $M$ tale che:
    \begin{itemize}[label = --]
        \item $M$ definisce una \textbf{riduzione} da $\mathcal{L}'$ a $\mathcal{L}$
        \item $M$ ha complessità compatibile (inferiore) a quella necessaria per riconoscere i linguaggi di $\mathbb{C}$
    \end{itemize}
\end{definition}}}\smallskip\\


%NP COMPLETEZZA

\section{NP-completezza}
Abbiamo visto che per definire problemi completi e difficili servono necessariamente le riduzioni. Per i problemi NP (\ref{np}) utilizzeremo \textbf{riduzioni polinomiali},
poichè una macchina deterministica con complessità temporale $\mathcal{O}(n^k)$ può essere simulata da una macchina non deterministica con la stessa complessità.\\
Una funzione $f: \Sigma^* \to \Gamma^*$ è calcolabile in tempo polinomiale se esiste una macchina di Turing $M$ tale che $M$ calcola $f$ e $M$ ha complessità temporale $\mathcal{O}(n^k)$
per qualche $k \in \mathbb{N}$.


\subsection{Riduzioni efficienti}
Le riduzioni efficienti le chiamiamo \textbf{riduzioni di Karp}, vediamo la definizione:\smallskip\\
\colorbox{SALZO}{\parbox{\dimexpr\linewidth-2\fboxsep}{\begin{definition}
    Siano $A$ e $B$ due linguaggi sugli alfabeti $\Sigma_A^*$ e $\Sigma_B^*$, rispettivamente. \\Diremo che 
    $\bm{A}$ \textbf{è Karp-riducibile in tempo polinomiale a B} (denotato con $A \le_{m}^{p}B$) se esiste una
    riduzione 
    \[f:\Sigma_A^* \to \Sigma_A^*\]
    da $A$ a $B$ tale che $f$ è calcolabile in tempo polinomiale  
\end{definition}}}\smallskip\\
La funzione $f$ viene chiamata \textbf{Karp-riduzione da} $\bm{A}$ in $\bm{B}$

%LINGUAGGI DIFFICILI E COMPLETI PER NP
\newpage
\subsection{Linguaggi difficili e completi per NP}
Ora che abbiamo introdotto le riduzioni efficienti, possiamo utilizzarle per definire linguaggi completi e difficili in NP. \smallskip\\
\colorbox{SALZO}{\parbox{\dimexpr\linewidth-2\fboxsep}{\begin{definition}
    Un linguaggio $\mathcal{L}$ si dice \textbf{NP-hard} sotto le riduzioni polinomiali se per ogni linguaggio $\mathcal{L}'\in NP$ è tale che  
    \[\mathcal{L}' \le^{p}_{m} \mathcal{L} \quad(\text{riduzione di efficiente Karp}) \]
\end{definition}}}\smallskip\\
In altre parole significa che il linguaggio $\mathcal{L}$ è almeno tanto difficile quanto qualsiasi problema della classe NP.\smallskip\\
\colorbox{SALZO}{\parbox{\dimexpr\linewidth-2\fboxsep}{\begin{definition}
    Un linguaggio $\mathcal{L}$ si dice \textbf{NP-complete} sotto le riduzioni in tempo polinomiale se $\mathcal{L} \in NP$ e $\mathcal{L} $ è NP-hard 
    sotto le riduzioni in tempo polinomiale 
\end{definition}}}\smallskip\\
In altre parole un linguaggio NP-completo è un linguaggio che appartiene alla classe NP ed è almeno tanto difficile quanto tutti gli altri problemi di NP.\smallskip\\
Per dare queste due definizioni abbiamo semplicemente applicato le definizioni generali viste nella sezione di introduzione del capitolo.\\
Successivamente vedremo il teorema di Cook e Levin, che si occupa di dimostrare che il linguaggio dei palindromi \textsc{sat} è NP-completo. 



\section{Teorema di Cook-Levin}
Per enunciare e dimostrare il teorema di Cook-Levin dobbiamo definire il linguaggio \textsc{sat}:\smallskip\\
\colorbox{SALZO}{\parbox{\dimexpr\linewidth-2\fboxsep}{\begin{definition}
    \textsc{sat} è il linguaggio di stringhe definito come segue:
    \[    \big\{    e(\varphi)  \ | \ \varphi \text{ è una formula proposizionale soddisfacibile su V}  \big\}      \]
    dove V è un insieme di formule proposizionali.
\end{definition}}}\smallskip\\
Il \textbf{Teorema di Cook-Levin} dimostra che il problema
della soddisfacibilità proposizionale ($SAT$) è
\textbf{NP-completo}.Fissiamo un encoding $e$ per le formule proposizionali,
utilizzando un alfabeto finito $\Sigma_{SAT}$.
L'encoding deve essere ragionevole, cioè deve permettere
di rappresentare e manipolare le formule in modo efficiente.\\
Enunciamo il teorema: \smallskip\\
\colorbox{SALZO}{\parbox{\dimexpr\linewidth-2\fboxsep}{
\begin{theorem}[Cook-Levin]
Il problema $SAT$ è \textbf{NP-hard} e, poiché
$SAT\in NP$, è anche \textbf{NP-completo}.
\end{theorem}
}}
Attenzione, la dimostrazione di questo teorema è abbastanza lunga e articolata poichè non possiamo appoggiarci 
a nessun risultato/linguaggi noti, mentre le prossime dimostrazioni saranno più semplici. 
\begin{proof}
\vspace{0.1cm}    
Per dimostrare che $SAT$ è NP-hard dobbiamo provare che
\textbf{ogni linguaggio appartenente a NP può essere
ridotto a SAT in tempo polinomiale}.
Consideriamo quindi un linguaggio arbitrario
$L\in NP$, definito su un alfabeto $\Sigma_L$.
Vogliamo costruire una funzione
\[
f_L:\Sigma_L^*\rightarrow\Sigma_{SAT}^*
\]
calcolabile in tempo polinomiale, tale che:
\[
\boxed{
x\in L\iff f_L(x)\in SAT
}
\]
Poiché $L\in NP$, per definizione esiste una macchina
di Turing non deterministica $M_L$ che decide $L$
in tempo polinomiale.
Quindi esiste $k\in\mathbb N$ tale che:
\[
t_{M_L}(n)=O(n^k).
\]
L'idea della dimostrazione consiste nel costruire,
per ogni input $x$, una formula proposizionale
$\varphi_x$ che rappresenti una possibile
esecuzione accettante di $M_L$ su $x$.\\
La formula deve soddisfare la seguente proprietà:
\[
\boxed{
M_L\text{ accetta }x
\iff
\varphi_x\text{ è soddisfacibile}
}
\]
Infatti:
\begin{itemize}
    \item Se esiste un'esecuzione accettante di $M_L$
    su $x$, deve esistere un assegnamento di verità
    che soddisfa $\varphi_x$.

    \item Se $\varphi_x$ è soddisfacibile, il suo
    assegnamento deve rappresentare un'esecuzione
    valida e accettante di $M_L$ su $x$.
\end{itemize}
Infine, dimostreremo che $\varphi_x$ può essere
costruita in tempo polinomiale rispetto alla
lunghezza dell'input $x$.\\
Poiché $M_L$ decide $L$ in tempo polinomiale,
esiste un polinomio $p(n)$ tale che:
\[
p(n)>\max\{n,t_{M_L}(n)\}.
\]
Il polinomio $p(n)$ dipende soltanto dalla macchina
$M_L$ e non dal particolare input.\\
Poniamo:
\[
n=|x|,\qquad T=p(n).
\]
Possiamo quindi stabilire che ogni esecuzione di
$M_L$ su un input di lunghezza $n$ utilizza
al più $T$ configurazioni.
Inoltre, una macchina di Turing può raggiungere
al massimo una nuova cella di nastro per transizione.
Di conseguenza, se l'esecuzione dura al più $T$
configurazioni, il numero di celle di nastro
necessarie è anch'esso limitato da $T$.
Pertanto possiamo rappresentare ogni esecuzione
considerando:
\begin{itemize}
    \item Al più $T$ istanti temporali.
    \item Al più $T$ celle del nastro.
    \item Un numero finito di stati della macchina.
\end{itemize}
Poiché la macchina termina in tempo polinomiale,
non può visitare un numero di celle superiore
a un limite polinomiale.
Questo ci permette di rappresentare un'intera
esecuzione mediante una tabella di dimensione
polinomiale rispetto all'input.
Consideriamo la macchina:
\[
M_L=
\langle
\Sigma,\Gamma,Q,\delta,q_0,q_{yes},q_{no}
\rangle.
\]
Una configurazione di $M_L$ è determinata da
tre informazioni:
\begin{enumerate}
    \item Lo \textbf{stato corrente} della macchina.
    \item La \textbf{posizione della testina} sul nastro.
    \item Il \textbf{contenuto delle celle} del nastro.
\end{enumerate}
Possiamo rappresentare un'esecuzione come una tabella in cui:
\begin{itemize}
    \item Ogni riga rappresenta una configurazione,
    ossia un istante temporale.

    \item Le colonne rappresentano lo stato della
    macchina, la posizione della testina e le $T$
    celle del nastro.
\end{itemize}
La tabella possiede dunque al più $T$ righe e
$T+2$ colonne.\medskip\\
\textit{Esempio.}
Supponiamo che una macchina abbia cinque celle
di nastro e che l'input sia $aa$.
Una possibile esecuzione può essere rappresentata
dalla seguente tabella:
\[
\begin{array}{c|c|ccccc}
\text{Stato} & \text{Testina}
& p_1 & p_2 & p_3 & p_4 & p_5\\
\hline
q_0     & 1 & a & a & \sqcup & \sqcup & \sqcup\\
q_0     & 2 & b & a & \sqcup & \sqcup & \sqcup\\
q_0     & 3 & b & b & \sqcup & \sqcup & \sqcup\\
q_1     & 3 & b & b & a & \sqcup & \sqcup\\
q_{yes} & 3 & b & b & a & \sqcup & \sqcup
\end{array}
\]
Ogni riga descrive completamente la situazione
della macchina in un determinato istante.
Il nostro obiettivo è costruire una formula
proposizionale che rappresenti una tabella di
questo tipo e garantisca che essa descriva
un'esecuzione valida.\\
Per rappresentare le configurazioni introduciamo
tre famiglie di variabili proposizionali.
Definiamo:
\[
H_j^i
\qquad i,j=1,\ldots,T.
\]
La variabile $H_j^i$ è vera se e solo se,
all'istante $i$, la testina si trova nella
posizione $j$ del nastro.
Formalmente:
\[
H_j^i=1
\iff
\text{testina in posizione }j
\text{ all'istante }i.
\]
Definiamo:
\[
Q_c^i
\qquad i=1,\ldots,T,\quad c\in Q.
\]
La variabile $Q_c^i$ è vera se e solo se
la macchina si trova nello stato $c$
all'istante $i$.
Formalmente:
\[
Q_c^i=1
\iff
\text{stato corrente }c
\text{ all'istante }i.
\]
Definiamo:
\[
V_g^{i,j}
\qquad
i,j=1,\ldots,T,\quad g\in\Gamma.
\]
La variabile $V_g^{i,j}$ è vera se e solo se,
all'istante $i$, la cella $j$ contiene
il simbolo $g$.
Formalmente:
\[
V_g^{i,j}=1
\iff
\text{la cella }j\text{ contiene }g
\text{ all'istante }i.
\]
\textit{Esempio.}
Consideriamo la prima configurazione della tabella:
\[
(q_0,1,aa\sqcup\sqcup\sqcup).
\]
Essa viene rappresentata assegnando valore vero
alle variabili:
\[
H_1^1\quad Q_{q_0}^1\quad
V_a^{1,1}\quad V_a^{1,2} \quad
V_{\sqcup}^{1,3}\quad
V_{\sqcup}^{1,4}\quad
V_{\sqcup}^{1,5}
\]
Le variabili che descrivono posizioni della testina,
stati o simboli diversi da quelli effettivi
assumono invece valore falso.
In questo modo, un assegnamento di verità
può rappresentare una possibile esecuzione.
Tuttavia, non tutti gli assegnamenti descrivono
esecuzioni valide: dobbiamo imporre dei vincoli.\smallskip\\
Costruiamo ora delle formule proposizionali
che garantiscano la correttezza di ogni
configurazione rappresentata.
In particolare, dobbiamo imporre tre condizioni:
\begin{enumerate}
    \item La testina occupa esattamente una posizione.
    \item La macchina si trova esattamente in uno stato.
    \item Ogni cella contiene esattamente un simbolo.
\end{enumerate}
Per ogni istante $i$, almeno una delle variabili
$H_1^i,\ldots,H_T^i$ deve essere vera:
\[
\bigvee_{j=1}^{T}H_j^i.
\]
Inoltre, due posizioni differenti non possono
essere contemporaneamente occupate:
\[
\bigwedge_{1\le j<k\le T}
(\neg H_j^i\lor\neg H_k^i).
\]
Combinando le due condizioni otteniamo:
\[
\boxed{
\varphi_H=
\bigwedge_{i=1}^{T}
\left[
\left(\bigvee_{j=1}^{T}H_j^i\right)
\quad
\land
\quad
\left(
\bigwedge_{1\le j<k\le T}
(\neg H_j^i\lor\neg H_k^i)
\right)
\right]
}
\]
La prima parte impone che la testina si trovi
\textbf{in almeno una posizione}; la seconda
impone che si trovi \textbf{in al più una posizione}.
Insieme garantiscono che occupi
\textbf{esattamente una posizione}.\smallskip\\
Analogamente, a ogni istante la macchina deve
trovarsi esattamente in uno stato dell'insieme $Q$.
Imponiamo:
\[
\boxed{
\varphi_Q=
\bigwedge_{i=1}^{T}
\left[
\left(\bigvee_{c\in Q}Q_c^i\right)
\quad
\land
\quad
\left(
\bigwedge_{\substack{c,d\in Q\\c\ne d}}
(\neg Q_c^i\lor\neg Q_d^i)
\right)
\right]
}
\]
La prima parte garantisce che almeno uno stato
sia attivo, mentre la seconda impedisce
che due stati distinti siano attivi
contemporaneamente.
Ogni cella del nastro deve contenere esattamente
un simbolo dell'alfabeto $\Gamma$.
Definiamo:
\[
\boxed{
\begin{aligned}
\varphi_V=
\bigwedge_{i=1}^{T}
\bigwedge_{j=1}^{T}
\Bigg[
\left(\bigvee_{g\in\Gamma}V_g^{i,j}\right)\quad
\land
\quad
\left(
\bigwedge_{\substack{g,h\in\Gamma\\g\ne h}}
(\neg V_g^{i,j}\lor\neg V_h^{i,j})
\right)
\Bigg]
\end{aligned}
}
\]
Questa formula garantisce che ogni cella
contenga un simbolo e non possa contenerne
due diversi nello stesso istante. Definiamo:
\[
\boxed{
\varphi_{\text{cons}}=
\varphi_H\land\varphi_Q\land\varphi_V
}
\]
Un assegnamento che soddisfa
$\varphi_{\text{cons}}$ rappresenta una sequenza
di configurazioni strutturalmente corrette.
Tuttavia, questo non garantisce ancora che
le configurazioni consecutive rispettino
le transizioni di $M_L$.\smallskip\\
Dobbiamo ora garantire che ogni configurazione
sia ottenuta dalla precedente applicando
una transizione valida della macchina.
Supponiamo che, nello stato $q$ e leggendo
il simbolo $a$, la funzione di transizione
non deterministica restituisca:
\[
\delta(q,a)=
\{
(q_1,b_1,d_1),\ldots,(q_s,b_s,d_s)
\}
\]
Questo significa che la macchina può scegliere
una delle $s$ transizioni disponibili.
Ogni transizione specifica:
\begin{itemize}
    \item Il nuovo stato $q_m$.
    \item Il simbolo $b_m$ da scrivere.
    \item La direzione $d_m$ dello spostamento.
\end{itemize}
Indichiamo con $j_m$ la posizione raggiunta
dalla testina applicando la transizione $m$:
\[
j_m=
\begin{cases}
j, & d_m=-,\\
j+1, & d_m=\rightarrow,\\
j-1, & d_m=\leftarrow,\ j>1,\\
1, & d_m=\leftarrow,\ j=1
\end{cases}
\]
Assumiamo, come nella costruzione delle slide,
che un movimento verso sinistra dalla prima cella
lasci la testina nella posizione 1.
Se all'istante $i$:
\begin{itemize}
    \item La testina è nella posizione $j$.
    \item La macchina è nello stato $q$.
    \item La cella $j$ contiene il simbolo $a$.
\end{itemize}
allora all'istante successivo deve verificarsi
una delle transizioni appartenenti a $\delta(q,a)$.
Questa condizione si rappresenta con:
\[
\begin{aligned}
H_j^i\land Q_q^i\land V_a^{i,j}
\quad\Rightarrow\quad
\bigvee_{m=1}^{s}
\left(
H_{j_m}^{i+1}
\land Q_{q_m}^{i+1}
\land V_{b_m}^{i+1,j}
\right)
\end{aligned}
\]
La parte a sinistra descrive la situazione
corrente della macchina.
La disgiunzione a destra rappresenta
\textbf{tutte le possibili transizioni valide}.
È sufficiente che una delle alternative
sia soddisfatta, perché una macchina
non deterministica deve scegliere
una sola transizione tra quelle disponibili.
Ripetiamo il vincolo per tutti gli istanti,
le posizioni, gli stati e i simboli:
\[
\boxed{
\begin{aligned}
\varphi_\delta=
\bigwedge_{i=1}^{T-1}
\bigwedge_{j=1}^{T}
\bigwedge_{q\in Q}
\bigwedge_{a\in\Gamma}
\Bigg[
&(H_j^i\land Q_q^i\land V_a^{i,j})
&\Rightarrow
\bigvee_{(q',b,d)\in\delta(q,a)}
\left(
H_{\mu(j,d)}^{i+1}
\land Q_{q'}^{i+1}
\land V_b^{i+1,j}
\right)
\Bigg]
\end{aligned}
}
\]
Qui $\mu(j,d)$ indica la nuova posizione
della testina dopo il movimento $d$.
Le transizioni che supererebbero il limite
delle $T$ celle vengono escluse dalla
rappresentazione: nessuna esecuzione valida
entro il limite temporale scelto ne ha bisogno.\smallskip\\
La formula precedente impone il comportamento
della cella attualmente letta dalla testina.
Dobbiamo però garantire che tutte le altre
celle rimangano invariate.
Se all'istante $i$ la testina non si trova
nella posizione $j$, il simbolo contenuto
nella cella $j$ deve essere lo stesso
all'istante $i+1$.
Formalmente:
\[
\boxed{
\begin{aligned}
\varphi_{\text{inv}}=
\bigwedge_{i=1}^{T-1}
\bigwedge_{j=1}^{T}
\bigwedge_{a\in\Gamma}
\Big[
(\neg H_j^i\land V_a^{i,j})
\Rightarrow V_a^{i+1,j}
\Big]
\end{aligned}
}
\]
Una transizione della macchina può modificare
soltanto la cella sulla quale si trova la
testina. Tutte le altre celle devono
conservare il proprio contenuto.\smallskip\\
Senza perdita di generalità, assumiamo che
gli stati $q_{yes}$ e $q_{no}$ abbiano
transizioni che mantengono invariati
lo stato, la posizione della testina
e il contenuto del nastro.
In questo modo, quando la macchina termina,
possiamo ripetere la configurazione terminale
nelle righe successive della tabella,
fino a raggiungere esattamente $T$ righe.
Questa convenzione semplifica la costruzione
delle formule.
Supponiamo che:
\[
\delta(q_0,a)=
\{
(q_0,b,\rightarrow),
(q_{yes},\sqcup,\leftarrow),
(q_{no},a,-)
\}
\]
Consideriamo l'istante $i=2$ e la posizione
$j=2$ del nastro.
Se la macchina si trova nello stato $q_0$
e legge $a$ nella seconda cella, può:
\begin{enumerate}
    \item Scrivere $b$, spostarsi a destra
    ed entrare nello stato $q_0$.

    \item Scrivere $\sqcup$, spostarsi a sinistra
    ed entrare nello stato $q_{yes}$.

    \item Lasciare $a$, non spostarsi
    ed entrare nello stato $q_{no}$.
\end{enumerate}
La corrispondente formula è:
\[
\begin{aligned}
H_2^2\land Q_{q_0}^2\land V_a^{2,2}
\quad
\Rightarrow
\quad
\bigg[
(H_3^3\land Q_{q_0}^3\land V_b^{3,2})
\quad\lor\quad
(H_1^3\land Q_{q_{yes}}^3\land V_{\sqcup}^{3,2})
\quad\lor\quad
(H_2^3\land Q_{q_{no}}^3\land V_a^{3,2})
\bigg]
\end{aligned}
\]
La disgiunzione rappresenta le tre possibili
scelte non deterministiche.
L'assegnamento di verità seleziona una
di queste alternative, rappresentando
una particolare esecuzione della macchina.\smallskip\\
Dobbiamo ora imporre che la prima riga della
tabella rappresenti esattamente la
configurazione iniziale di $M_L$ su $x$.
Sia:
\[
x=c_1c_2\cdots c_n
\]
All'istante iniziale:
\begin{itemize}
    \item La macchina si trova nello stato $q_0$.
    \item La testina si trova nella posizione 1.
    \item Le prime $n$ celle contengono l'input.
    \item Le celle successive contengono il
    simbolo bianco $\sqcup$.
\end{itemize}
Definiamo quindi:
\[
\boxed{
\begin{aligned}
\varphi_{\text{init}}={}&
Q_{q_0}^{1}\quad\land\quad H_1^1
\quad
\land
\quad
\bigwedge_{j=1}^{n}V_{c_j}^{1,j}
\quad
\land
\quad
\bigwedge_{j=n+1}^{T}V_{\sqcup}^{1,j}
\end{aligned}
}
\]
Per l'input vuoto, tutte le celle iniziali
sono bianche.\smallskip\\
\textit{Esempio.}
Se $x=aa$ e $T=5$, otteniamo:
\[
\begin{aligned}
\varphi_{\text{init}}={}&
Q_{q_0}^{1} \quad\land\quad H_1^1
\quad\land \quad V_a^{1,1}
\quad\land\quad V_a^{1,2}
\quad
\land\quad V_{\sqcup}^{1,3}
\quad\land\quad V_{\sqcup}^{1,4}
\quad\land \quad V_{\sqcup}^{1,5}
\end{aligned}
\]
Questa formula descrive la configurazione
iniziale della macchina con input $aa$.\smallskip\\
Infine, vogliamo imporre che l'esecuzione
rappresentata dalla tabella sia accettante.
Per questo è sufficiente richiedere che
lo stato $q_{yes}$ venga raggiunto
in almeno uno degli istanti temporali.
Definiamo:
\[
\boxed{
\varphi_{\text{acc}}=
\bigvee_{i=1}^{T}Q_{q_{yes}}^i
}
\]
La formula è vera se esiste almeno un
istante nel quale la macchina si trova
nello stato accettante.
\textit{Esempio.}
Per $T=5$:
\[
\varphi_{\text{acc}}=
Q_{q_{yes}}^1\quad\lor\quad Q_{q_{yes}}^2
\quad\lor\quad Q_{q_{yes}}^3\quad\lor\quad Q_{q_{yes}}^4
\quad\lor\quad Q_{q_{yes}}^5
\]
È sufficiente che una di queste variabili
sia vera per soddisfare il vincolo.\smallskip\\
Abbiamo costruito diverse formule,
ciascuna delle quali impone una particolare
condizione sull'esecuzione.
Le combiniamo mediante una congiunzione:
\[
\boxed{
\begin{aligned}
\varphi_x={}
\varphi_{\text{cons}}
\quad\land\quad\varphi_\delta
\quad\land\quad\varphi_{\text{inv}}
\quad\land\quad\varphi_{\text{init}}
\quad\land\quad\varphi_{\text{acc}}
\end{aligned}
}
\]
Le componenti hanno i seguenti significati:
\begin{itemize}
    \item $\varphi_{\text{cons}}$:
    ogni configurazione è strutturalmente valida.

    \item $\varphi_\delta$:
    ogni passaggio tra configurazioni
    rispetta una transizione di $M_L$.

    \item $\varphi_{\text{inv}}$:
    le celle non lette dalla testina
    mantengono il proprio contenuto.

    \item $\varphi_{\text{init}}$:
    l'esecuzione parte dalla configurazione
    iniziale corretta per l'input $x$.

    \item $\varphi_{\text{acc}}$:
    l'esecuzione raggiunge uno stato accettante.
\end{itemize}
Pertanto, soddisfare $\varphi_x$ significa
poter descrivere una sequenza valida
di configurazioni che parte dall'input $x$
e raggiunge lo stato accettante.\smallskip\\
\colorbox{SALZO}{\parbox{\dimexpr\linewidth-2\fboxsep}{
\begin{theorem}[\textbf{Lemma 1}]
La formula $\varphi_x$ è soddisfacibile
se e solo se esiste un'esecuzione accettante
della macchina $M_L$ sull'input $x$.
\end{theorem}
}}\smallskip\\
Dimostriamo entrambe le implicazioni.\smallskip\\
$\bullet$ \textbf{Prima direzione}:
Supponiamo che $M_L$ abbia un'esecuzione
accettante sull'input $x$.
Poiché $M_L$ termina in tempo polinomiale,
l'esecuzione può essere rappresentata
con al più $T$ configurazioni e $T$ celle.
Costruiamo un assegnamento di verità alle
variabili proposizionali nel seguente modo:
\begin{itemize}[label = --]
    \item Rendiamo vera $H_j^i$ se la testina,
    all'istante $i$, si trova nella posizione $j$.

    \item Rendiamo vera $Q_c^i$ se lo stato
    corrente all'istante $i$ è $c$.

    \item Rendiamo vera $V_g^{i,j}$ se
    la cella $j$, all'istante $i$, contiene $g$.
\end{itemize}
Poiché l'esecuzione è valida:
\begin{itemize}[label = --]
    \item I vincoli di consistenza sono rispettati.
    \item Le transizioni corrispondono a quelle
    della macchina.
    \item Le celle non modificate restano invariate.
    \item La configurazione iniziale è corretta.
    \item Viene raggiunto lo stato accettante.
\end{itemize}
L'assegnamento soddisfa quindi tutti
i congiunti di $\varphi_x$.
Pertanto:
\[
M_L\text{ accetta }x
\Rightarrow
\varphi_x\text{ è soddisfacibile}.
\]
$\bullet$ \textbf{Seconda direzione}:
Supponiamo che $\varphi_x$ sia soddisfacibile.
Esiste quindi un assegnamento di verità
che rende veri tutti i vincoli della formula.
Dai vincoli di consistenza possiamo ricostruire
una sequenza di configurazioni ben definite.
Dai vincoli sulle transizioni sappiamo
che ogni configurazione successiva è ottenuta
applicando una transizione valida di $M_L$.
I vincoli di inizializzazione garantiscono
che la sequenza parta dalla configurazione
iniziale di $M_L$ sull'input $x$.
Infine, la condizione di accettazione garantisce
che almeno una configurazione raggiunga
lo stato $q_{yes}$.
Di conseguenza, la sequenza rappresenta
un'esecuzione valida e accettante di $M_L$
sull'input $x$.
Pertanto:
\[
\varphi_x\text{ è soddisfacibile}
\Rightarrow
M_L\text{ accetta }x.
\]
Avendo dimostrato entrambe le implicazioni:
\[
\boxed{
M_L\text{ accetta }x
\iff
\varphi_x\text{ è soddisfacibile}.
}
\]
Per completare la dimostrazione dobbiamo
mostrare che $\varphi_x$ può essere costruita
in \textbf{tempo polinomiale rispetto a $|x|$}.\smallskip\\
\colorbox{SALZO}{\parbox{\dimexpr\linewidth-2\fboxsep}{
\begin{theorem}[\textbf{Lemma 2}]
La formula proposizionale $\varphi_x$
può essere costruita in tempo polinomiale
rispetto alla lunghezza dell'input $x$.
\end{theorem}
}}\smallskip\\
Dimostriamo questo lemma ricordando che:
\[
T=p(n),\qquad n=|x|
\]
dove $p(n)$ è un polinomio.
Calcoliamo innanzitutto quante variabili
proposizionali sono necessarie.\\
Per ogni istante temporale abbiamo:
\begin{itemize}
    \item $T$ variabili per la posizione
    della testina.

    \item $|Q|$ variabili per lo stato corrente.

    \item $|\Gamma|\cdot T$ variabili
    per il contenuto del nastro.
\end{itemize}
Il totale per una configurazione è quindi:
\[
T+|Q|+|\Gamma|T
\]
Poiché consideriamo $T$ configurazioni,
il numero complessivo di variabili è:
\[
\boxed{
N_{\text{var}}=
T(T+|Q|+|\Gamma|T)
}
\]
Sviluppando:
\[
N_{\text{var}}=
(1+|\Gamma|)T^2+|Q|T
\]
Poiché $M_L$ è fissata, gli insiemi
$Q$ e $\Gamma$ hanno dimensione costante.
Quindi:
\[
\boxed{N_{\text{var}}=O(T^2)}
\]
Essendo $T=p(n)$ un polinomio,
anche il numero delle variabili è
polinomiale rispetto a $n$.\smallskip\\
Consideriamo ora le formule costruite
precedentemente:
\begin{center}
\begin{tabular}{|l|c|}
\hline
\textbf{Tipologia di vincolo}
& \textbf{Quantità} \\
\hline
Consistenza della testina & $O(T^3)$ \\
\hline
Consistenza dello stato & $O(T)$ \\
\hline
Consistenza dei simboli & $O(T^2)$ \\
\hline
Transizioni & $O(T^2)$ \\
\hline
Celle invariate & $O(T^2)$ \\
\hline
Inizializzazione e accettazione & $O(T)$ \\
\hline
\end{tabular}
\end{center}
La quantità maggiore è quella dei vincoli
di consistenza sulla posizione della testina,
pari a $O(T^3)$.
Questo accade perché, per ciascuno dei $T$
istanti, dobbiamo impedire che due delle $T$
posizioni siano occupate contemporaneamente.
Il numero delle coppie di posizioni è:
\[
\binom{T}{2}=\frac{T(T-1)}{2}=O(T^2).
\]
Ripetendo il controllo per tutti gli istanti:
\[
T\cdot O(T^2)=O(T^3)
\]
Le altre tipologie di vincoli richiedono
un numero inferiore di formule elementari.
Poiché la macchina $M_L$ è fissata,
le dimensioni di $Q$, $\Gamma$ e della
funzione di transizione sono costanti
rispetto alla lunghezza dell'input.
La formula completa ha quindi dimensione
polinomiale in $T$.
Essendo:
\[
T=p(n)
\]
con $p(n)$ polinomiale, anche la dimensione
di $\varphi_x$ è polinomiale rispetto a $n$.
Inoltre, i vincoli possono essere generati
sistematicamente enumerando gli istanti,
le posizioni, gli stati e i simboli necessari.
Pertanto la costruzione non richiede soltanto
uno spazio polinomiale per la formula,
ma anche un \textbf{tempo polinomiale}
per produrla.
Concludiamo:
\[
\boxed{
\varphi_x\text{ è costruibile in tempo polinomiale}.
}
\]
Abbiamo dimostrato che, per ogni input $x$,
è possibile costruire una formula $\varphi_x$
con le seguenti proprietà:
\begin{enumerate}
    \item $\varphi_x$ è soddisfacibile
    se e solo se $M_L$ accetta $x$.

    \item $\varphi_x$ può essere costruita
    in tempo polinomiale rispetto a $|x|$.
\end{enumerate}
Definiamo quindi la funzione:
\[
\boxed{f_L(x)=e(\varphi_x).}
\]
Questa funzione associa a ogni input $x$
la codifica della formula costruita.\smallskip\\
\colorbox{SALZO}{\parbox{\dimexpr\linewidth-2\fboxsep}{
\begin{theorem}[Lemma 3]
La funzione $f_L$ è una riduzione
molti-a-uno polinomiale da $L$ a $SAT$.
\end{theorem}
}}\smallskip\\
Per il Lemma 2, la funzione $f_L$
è calcolabile in tempo polinomiale.
Inoltre, per il Lemma 1:
\[
\begin{aligned}
x\in L
&\iff M_L\text{ accetta }x\\
&\iff \varphi_x\text{ è soddisfacibile}\\
&\iff e(\varphi_x)\in SAT\\
&\iff f_L(x)\in SAT
\end{aligned}
\]
Pertanto:
\[
\boxed{L\le_p SAT}
\]
Abbiamo considerato un linguaggio arbitrario
$L\in NP$ e dimostrato che:
\[
L\le_p SAT
\]
Poiché la costruzione vale per
\textbf{qualsiasi linguaggio appartenente a NP},
possiamo concludere che:
\[
\boxed{
\forall L\in NP,\quad L\le_p SAT.
}
\]
Per definizione, questo significa che
$SAT$ è NP-hard.\\
Inoltre, sappiamo che:
\[
SAT\in NP
\]
poiché una macchina non deterministica
può scegliere un assegnamento di verità
e verificare in tempo polinomiale
se esso soddisfa la formula.
Dunque $SAT$ è sia NP-hard sia appartenente
a NP.\smallskip\\
Possiamo finalmente concludere:
\[
\boxed{
SAT\text{ è NP-completo}
}
\]
\end{proof}

\noindent Il Teorema di Cook-Levin dimostra che
\textbf{qualsiasi computazione non deterministica
di durata polinomiale può essere rappresentata
tramite una formula proposizionale}.
In particolare:
\begin{itemize}
    \item Le variabili descrivono le configurazioni
    della macchina.

    \item I vincoli garantiscono che le configurazioni
    rappresentino un'esecuzione valida.

    \item La condizione di accettazione impone
    che l'esecuzione raggiunga $q_{yes}$.

    \item La formula è soddisfacibile esattamente
    quando esiste un'esecuzione accettante.

    \item La costruzione richiede soltanto
    tempo polinomiale.
\end{itemize}
Di conseguenza, qualsiasi problema in NP
può essere trasformato in un'istanza di SAT
senza introdurre un aumento più che polinomiale
della dimensione dell'input.
Questo rende SAT uno dei problemi fondamentali
della teoria della complessità computazionale.


